\documentclass[12pt]{article}
\usepackage{amsmath}
\usepackage{mathrsfs}
\usepackage{eucal}
\usepackage{amssymb}
\usepackage[french]{babel}
\usepackage[latin1]{inputenc}
\usepackage[T1]{fontenc}
\usepackage[all]{xy}
%\usepackage{graphicx}
%\usepackage{epsf}

\begin{document}

\hfuzz 8 pt
\parindent = 0 cm

\centerline{ \textsc{ \large Universit� de Sherbrooke}}
\centerline{\textsc{ \normalsize D�partement de math�matiques}}
\centerline{\textsc{ \normalsize �t� 2005}}

\bigskip

\centerline{\large\bf MAT417 -- M�thodes num�riques en alg�bre
lin�aire}

\bigskip
\centerline{\Large\bf Laboratoire 9}
\bigskip


\begin{description}

\item[1)] Consid�rons la matrice

$$A = \left(
\begin{array}{rrrr}
0,7 & -1,7 & -2,9 & 0,5 \\
-1,7 & 0,7 & -0,5 & 2,9 \\
-2,4 & 2,4 & 2,2 & -2,2 \\
-2,4 & 2,4 & -0,2 & 0,2
\end{array}
\right)
$$

    \begin{description}

    \item[a)] D�terminer $\lambda_1$ et $x_1$ en appliquant la
    m�thode des puissances � partir du vecteur initial $z_0 =
    (1,-1,0,0)^t$.

    \item[b)] Calculer $\lambda_4$ par la m�thode des puissances
    inverses appliqu�e � partir du vecteur $z_0 = (0,0,1,1)^t$,
    puis calculer $x_4$.

    \item[c)] D�terminer la valeur de $\lambda_2$ en appliquant la
    m�thode des puissances inverses � la matrice $A-2,5 1\!{\rm
    l}$ � partir de $z_0 =
    (\frac{1}{2},\frac{1}{2},\frac{1}{2},\frac{1}{2})^t$.

    \item[d)] D�terminer la derni�re valeur propre � l'aide d'une
    m�thode de votre choix.

    \end{description}


\item[2)] Consid�rons la matrice

$$A = \frac{1}{9} \left(
\begin{array}{rrr}
5 & 2 & 2 \\
2 & -1 & 8 \\
2 & 8 & -1
\end{array}
\right)
$$

D�terminer la valeur propre dominante de $A$ et un vecteur propre
associ� par la m�thode des puissances appliqu�e � partir du
vecteur $z_0 = (1,2,0)^t$.

\end{description}

\end{document}
